\chapter*{Abstract}\addcontentsline{toc}{chapter}{Abstract}

% \begin{explanation}
%     An abstract is a concise summary of your manuscript and must make sense by itself. It is probably the first part that a reader reads. It must describe the work and set accurate expectations for the reader. It is usually around 200 to 300 words and provides the following information: the context and reason for writing, the problem statement and scope, the used method(s) and/or performed experiments, the obtained results and the conclusions and implications.
    
%     An abstract in English is required for a master thesis in Engineering Technology and BRUFACE.
% \end{explanation}

% The text of the abstract in this template is added with the \input command below
% To remove it, put a % in front of it,
% as follows: %This is a template to make a thesis and/or slides in \href{https://www.latex-project.org}{\LaTeX}, free software to compile high-quality scientific documents. The template contains the typical parts such as a title page (for VUB, BRUFACE and ULB), acknowledgements, abstract, table of contents, a nomenclature, the typical chapters, appendices and a bibliography. In the appendices you will find practical information about formatting and examples of useful packages.

This template is too large to compile with a free Overleaf plan. Consult section A.1.4 in appendix A for possible workarounds. % a hard link to A.1.4 is added on purpose, because \cref will not work when appendix A is not included, as in myThesis.tex

A \href{https://vub.cloud.panopto.eu/Panopto/Pages/Sessions/List.aspx?folderID=eb5cbed1-a7f8-4d08-bca3-ad1e00db8d7a}{set of videos} is available that demonstrate the use of this template.

This template comes with four main files:
\begin{itemize}
    \item \texttt{mainReport}: to compile a report and all the appendices of this template
    \item \texttt{mainSlides}: to compile slides. Only selected content is put on the slides though. To see all content, you must compile \texttt{mainReport} and/or inspect the source code.
    \item \texttt{PhD}: to compile a PhD without the appendices of this template
    \item \texttt{Thesis}: to compile a thesis without the appendices of this template
\end{itemize}

The template uses the class \href{https://ctan.org/pkg/memoir?lang=en}{memoir}. It is a class for typesetting mathematical works as books, reports and articles.

This template is written in English. If you want to write in Dutch, you must select Dutch as the spell checking language in Overleaf via the Menu-button in the upper left corner and replace \texttt{english} by \texttt{dutch} as option in the \cs{documentclass} command.
%This is a template to make a thesis and/or slides in \href{https://www.latex-project.org}{\LaTeX}, free software to compile high-quality scientific documents. The template contains the typical parts such as a title page (for VUB, BRUFACE and ULB), acknowledgements, abstract, table of contents, a nomenclature, the typical chapters, appendices and a bibliography. In the appendices you will find practical information about formatting and examples of useful packages.

This template is too large to compile with a free Overleaf plan. Consult section A.1.4 in appendix A for possible workarounds. % a hard link to A.1.4 is added on purpose, because \cref will not work when appendix A is not included, as in myThesis.tex

A \href{https://vub.cloud.panopto.eu/Panopto/Pages/Sessions/List.aspx?folderID=eb5cbed1-a7f8-4d08-bca3-ad1e00db8d7a}{set of videos} is available that demonstrate the use of this template.

This template comes with four main files:
\begin{itemize}
    \item \texttt{mainReport}: to compile a report and all the appendices of this template
    \item \texttt{mainSlides}: to compile slides. Only selected content is put on the slides though. To see all content, you must compile \texttt{mainReport} and/or inspect the source code.
    \item \texttt{PhD}: to compile a PhD without the appendices of this template
    \item \texttt{Thesis}: to compile a thesis without the appendices of this template
\end{itemize}

The template uses the class \href{https://ctan.org/pkg/memoir?lang=en}{memoir}. It is a class for typesetting mathematical works as books, reports and articles.

This template is written in English. If you want to write in Dutch, you must select Dutch as the spell checking language in Overleaf via the Menu-button in the upper left corner and replace \texttt{english} by \texttt{dutch} as option in the \cs{documentclass} command. % 

% write your own English abstract here
Belgian residential energy demand remains dominated by space heating, but climate change and electrification alter both the magnitude and timing of future demand. This thesis develops a stochastic, physics-informed bottom-up model of Belgian dwellings to analyse how residential final energy, electricity demand, peak load, PV self-consumption, emissions, and costs respond to climate forcing and counterfactual technology stocks. The model combines stock-weighted dwelling archetypes, stochastic occupancy and end-use behaviour, a reduced-order thermal balance, temperature-dependent heat-pump COPs, PV and EV modules, and a scenario-tree design that varies four dimensions: climate window, RCP pathway, technology case, and stochastic realisation.
\newline The validation strategy separates calibration from external validation. Macro-level end-use totals and PV performance are reproduced strongly, synthetic diversity behaviour is plausible, and EV timing is reproduced well after magnitude sensitivity. The Fluvius smart-meter comparison does not converge at the household aggregate load-shape level, so aggregate profile and peak-timing results are interpreted as directional rather than as externally calibrated forecasts. The executed scenario evidence covers 180 completed leaves of the planned 2800-leaf tree and is therefore used as conditional scenario evidence rather than as a fully sampled probabilistic ensemble.
\newline Within that evidence base, climate warming mainly reduces useful heating pressure while increasing cooling pressure. In the long-term RCP\,8.5 case, HDD$_{18}$ falls by about \SI{31}{\percent}, CDD$_{22}$ rises from about $14.3$ to $109.3\,^{\circ}\mathrm{C}\cdot\mathrm{d}$ per year, and useful space-heating demand falls by about \SI{33}{\percent}. By contrast, peak electricity stress is driven primarily by technology electrification: in the cold-design-year comparison, frozen stock remains near the historical peak at about \SI{31}{kW}, while moderate electrification raises the cohort peak to about \SI{106}{kW} and the high-electrification PV--EV case to about \SI{188}{kW}. The main contribution is therefore a traceable conditional modelling framework showing that future Belgian residential demand is not only a heating-demand problem, but a coupled electrification, coincidence, cooling-pressure, PV self-consumption, tariff, and emissions problem.

\bigskip
\textbf{keywords}: residential energy demand, bottom-up modelling, climate scenarios, electrification, \LaTeX, \TikZ


% Abstract in Dutch
\chapter*{Samenvatting}\addcontentsline{toc}{chapter}{Samenvatting}

De Belgische residentiële energievraag blijft gedomineerd door ruimteverwarming, maar klimaatverandering en elektrificatie veranderen zowel de omvang als het tijdsprofiel van toekomstige vraag. Deze thesis ontwikkelt een stochastisch, fysisch geïnformeerd bottom-upmodel van Belgische woningen om te analyseren hoe residentiële finale energie, elektriciteitsvraag, piekbelasting, PV-zelfverbruik, emissies en kosten reageren op klimaatforcering en contrafeitelijke technologiestocks. Het model combineert stockgewogen woningarchetypes, stochastische bezetting en eindgebruiksgedrag, een gereduceerde thermische balans, temperatuurafhankelijke COP’s van warmtepompen, PV- en EV-modules, en een scenarioboomontwerp dat vier dimensies varieert: klimaatvenster, RCP-pad, technologiecase en stochastische realisatie.
\newline De validatiestrategie maakt onderscheid tussen kalibratie en externe validatie. Macro-totalen en PV-prestaties worden sterk gereproduceerd, synthetisch diversiteitsgedrag is plausibel, en EV-timing wordt goed gereproduceerd na gevoeligheidsanalyse van de grootteorde. De Fluvius smart-metervergelijking convergeert niet op het niveau van geaggregeerde huishoudelijke load shape, waardoor geaggregeerde profiel- en piektimingresultaten als richtinggevend worden geïnterpreteerd in plaats van als extern gekalibreerde voorspellingen. De uitgevoerde scenario-evidentie omvat 180 voltooide bladeren van de geplande scenarioboom met 2800 bladeren en wordt daarom gebruikt als voorwaardelijke scenario-evidentie, niet als volledig bemonsterd probabilistisch ensemble.
\newline Binnen die evidentiebasis verlaagt klimaatopwarming vooral de nuttige verwarmingsdruk, terwijl de koeldruk toeneemt. In het langetermijngeval RCP 8.5 daalt HDD$_{18}$ met ongeveer \SI{31}{\percent}, stijgt CDD$_{22}$ van ongeveer $14.3$ naar $109.3\,^{\circ}\mathrm{C}\cdot\mathrm{d}$ per jaar, en daalt de nuttige ruimteverwarmingsvraag met ongeveer \SI{33}{\percent}. Daartegenover wordt elektrische piekbelasting vooral gedreven door technologische elektrificatie: in de vergelijking met het koude ontwerpjaar blijft de frozen stock dicht bij de historische piek van ongeveer \SI{31}{kW}, terwijl matige elektrificatie de cohortpiek verhoogt tot ongeveer \SI{106}{kW} en de hoge-elektrificatie PV--EV-case tot ongeveer \SI{188}{kW}. De belangrijkste bijdrage is daarom een traceerbaar voorwaardelijk modelleringskader dat aantoont dat toekomstige Belgische residentiële vraag niet alleen een verwarmingsvraagstuk is, maar een gekoppeld probleem van elektrificatie, gelijktijdigheid, koeldruk, PV-zelfverbruik, tarieven en emissies.

\bigskip
\textbf{trefwoorden}: residentiële energievraag, bottom-upmodellering, klimaatscenario’s, elektrificatie, \LaTeX, \TikZ


% Abstract in French
\chapter*{R\'esum\'e}\addcontentsline{toc}{chapter}{R\'esum\'e}

La demande énergétique résidentielle belge reste dominée par le chauffage des locaux, mais le changement climatique et l’électrification modifient à la fois l’ampleur et le profil temporel de la demande future. Ce mémoire développe un modèle bottom-up stochastique et physiquement informé des logements belges afin d’analyser la manière dont l’énergie finale résidentielle, la demande d’électricité, la pointe de charge, l’autoconsommation PV, les émissions et les coûts réagissent au forçage climatique et à des stocks technologiques contrefactuels. Le modèle combine des archétypes de logements pondérés par le stock, une occupation et des usages finaux stochastiques, un bilan thermique réduit, des COP de pompes à chaleur dépendants de la température, des modules PV et EV, ainsi qu’un arbre de scénarios faisant varier quatre dimensions : fenêtre climatique, trajectoire RCP, cas technologique et réalisation stochastique.
\newline La stratégie de validation distingue la calibration de la validation externe. Les totaux macro par usage final et la performance PV sont fortement reproduits, le comportement de diversité synthétique est plausible, et le profil temporel EV est bien reproduit après analyse de sensibilité sur la magnitude. La comparaison avec les compteurs intelligents de Fluvius ne converge pas au niveau de la forme de charge agrégée des ménages ; les résultats relatifs aux profils agrégés et au moment de la pointe sont donc interprétés comme directionnels plutôt que comme des prévisions calibrées externement. Les résultats de scénarios exécutés couvrent 180 feuilles terminées sur les 2800 prévues et sont donc utilisés comme évidence conditionnelle de scénario, et non comme un ensemble probabiliste entièrement échantillonné.
\newline Dans cette base d’évidence, le réchauffement climatique réduit principalement la pression de chauffage utile tout en augmentant la pression de refroidissement. Dans le cas à long terme RCP\,8.5, les HDD$_{18}$ diminuent d’environ \SI{31}{\percent}, les CDD$_{22}$ augmentent d’environ $14.3$ à $109.3\,^{\circ}\mathrm{C}\cdot\mathrm{d}$ par an, et la demande utile de chauffage des locaux diminue d’environ \SI{33}{\percent}. En revanche, la contrainte de pointe électrique est principalement déterminée par l’électrification technologique : dans la comparaison de l’année de conception froide, le stock figé reste proche de la pointe historique, à environ \SI{31}{kW}, tandis qu’une électrification modérée porte la pointe de cohorte à environ \SI{106}{kW} et que le cas PV--EV à forte électrification l’élève à environ \SI{188}{kW}. La contribution principale est donc un cadre de modélisation conditionnel et traçable montrant que la demande résidentielle belge future n’est pas seulement un problème de demande de chauffage, mais un problème couplé d’électrification, de coïncidence, de pression de refroidissement, d’autoconsommation PV, de tarifs et d’émissions.

\bigskip
\textbf{mots-clés}: demande énergétique résidentielle, modélisation bottom-up, scénarios climatiques, électrification, \LaTeX, \TikZ


\mode<all> % do not remove
